% ECONOMICS_PROBLEM: {"id": "D91-558", "title": "Scaling behavioral interventions", "jel_code": "D91", "source_id": "OP-0308"}
% !TeX program = xelatex
\documentclass[11pt,a4paper]{article}
\usepackage[margin=23mm]{geometry}
\usepackage{fontspec}
\setmainfont{Latin Modern Roman}
\usepackage{amsmath,amssymb,mathtools}
\usepackage{xcolor,enumitem,booktabs,longtable,array,xurl,hyperref}
\definecolor{muted}{HTML}{53636C}
\hypersetup{colorlinks=true,urlcolor=blue,linkcolor=blue}
\setlength{\parindent}{0pt}
\setlength{\parskip}{5pt}
\newcommand{\R}{\mathbb R}
\newcommand{\E}{\mathbb E}
\newcommand{\N}{\mathbb N}
\newcommand{\ind}{\mathbf 1}
\newcommand{\Var}{\operatorname{Var}}
\newcommand{\Cov}{\operatorname{Cov}}
\newcommand{\supp}{\operatorname{supp}}
\DeclareMathOperator*{\argmin}{arg\,min}
\DeclareMathOperator*{\argmax}{arg\,max}
\newcommand{\entrymeta}[3]{{\sffamily\small\color{muted}\textbf{\texttt{#1}}\quad|\quad #2\par #3\par}}
\newcommand{\Problemref}[1]{\texttt{#1}}
\newcommand{\JELref}[2]{\texttt{#2} (JEL \texttt{#1})}
\begin{document}
\section*{Scaling behavioral interventions}
\textbf{Problem ID:} \texttt{D91-558}\quad \textbf{JEL:} \texttt{D91}\par
% BEGIN ECONOMICS STATEMENT
\label{problem:OP-0308}
{\sffamily\small\color{muted}Original subject group: Behavioral and experimental economics\par}
\label{definitions:problem:30007590}
\textbf{Audit classification:} Published research agenda. DellaVigna–Linos March 2020 manuscript explicitly asks about scaling and leaves forecasting accuracy needing more trials. The causal saturation surface is editorial.
\subsection*{Definitions and precise research target}
\textbf{Model and research target.} Index interventions by \(j\) and implementation scale by \(s\), defined as the fraction of an eligible population reached. Let \(Y_{ij}(d,s)\) be participant \(i\)'s outcome under receipt \(d\in\{0,1\}\), allowing implementation capacity and spillovers to depend on scale. Define the direct treatment effect \(\tau_j(s)=\mathbb E[Y_{ij}(1,s)-Y_{ij}(0,s)]\). Let \(X_j\) contain preregistered intervention features and \(C_j(s)\) contain measured implementation features. The prediction target for a large-scale rollout is
\[\widehat\tau_j(s_H)=g\bigl(\widehat\tau_j(s_L),X_j,C_j(s_L),C_j(s_H)\bigr),\qquad R=\mathbb E[(\widehat\tau_j(s_H)-\tau_j(s_H))^2].\]
Estimate and validate \(g\) on interventions held out from training, using complete trial registries rather than only published significant effects. Separate differences caused by publication selection, population composition, intervention versions, and administrative constraints from the causal effect of changing scale itself. Observing academic and government trials at different scales does not identify \(\tau_j(s_H)-\tau_j(s_L)\) for a common intervention. Identification of that contrast requires randomized saturation, replicated implementation, or explicit transport assumptions. Assess both rank prediction and magnitude prediction, and include uncertainty from estimated trial effects in \(R\).

\emph{Definitions and maintained assumptions (editorial unless a source definition is named).} Let \(s\in[0,1]\) denote a specified treatment saturation within a defined interference cluster; exposure mapping must justify \(Y_{ij}(d,s)\) rather than dependence on every other individual’s assignment. \(\tau_j(s)\) is the direct treatment effect at saturation \(s\), not total rollout welfare or the outcome effect on untreated people. Randomized two-stage assignment first sets cluster saturation, then individual receipt, with support for \(d=0,1\) at each studied \(s\). Population composition, implementation version and implementer must be held fixed or included in the transport model. Define a target law \(\mathcal J\) over new interventions and implementations for risk \(R=\mathbb E_{j\sim\mathcal J}[(\widehat\tau_j(s_H)-\tau_j(s_H))^2]\); without \(\mathcal J\), prediction risk is unspecified. Training and test interventions are independent under a declared sampling scheme. Trial-effect estimation noise must be separated from latent forecast error. Complete registries reduce publication selection but do not randomize implementation scale.
\subsection*{Variants and assumption changes}
\begin{enumerate}
\item \textbf{Causal scale effect (editorial).} Randomize saturation for the same intervention, implementer and population. Estimate both direct effects \(\tau_j(s)\) and spillover contrasts among untreated units.
\item \textbf{Cross-institution forecast (source-related editorial).} Train on registered small and large trials with specified feature coding and target intervention law. Test predictions on entirely held-out interventions, reporting transport assumptions rather than treating academic--government differences as scale effects.
\end{enumerate}
\subsection*{References and source audit}
\begin{enumerate}
\item §1 p.1; end of §4 p.21 in March2020 manuscript. Supports the stated research area; exact displayed specification is editorial. \url{https://eml.berkeley.edu/~sdellavi/wp/NudgeToScale2020-03-20.pdf}
\end{enumerate}
\begin{enumerate}\raggedright
\item\hypertarget{definitions:source:30007590:1}{} Stefano DellaVigna; Elizabeth Linos. 2020. \emph{RCTs to Scale: Comprehensive Evidence from Two Nudge Units}. Author manuscript dated March 2020, §1, printed p. 1; end of §4 and start of §5, printed p. 21 (PDF pp. 2 and 22)..
\url{https://eml.berkeley.edu/~sdellavi/wp/NudgeToScale2020-03-20.pdf}.
\end{enumerate}

\subsection*{Common conventions: Source mathematical conventions}

These conventions apply to each entry unless explicitly replaced.

\textbf{Models and identification.} A primitive model \(m\) specifies all probability laws, feasible choices, information, equilibrium and measurement rules used by its target. The admissible class \(\mathcal M\) is fixed independently of outcomes. If \(P_O(m)\) is its observable probability law and \(\tau(m)\) the target, its identified set at observable law \(P\) is
\[
 \mathcal I_\tau(P)=\{\tau(m):m\in\mathcal M,\ P_O(m)=P\}.
\]
Point identification means that this set is a singleton. An empty set signals incompatibility of data and assumptions. Coordinate bounds need not describe the joint set. Bounds are sharp only if every retained target is attainable or arbitrarily approachable by compatible models. Finite bounds require explicit restrictions on unobserved quantities; unrestricted confounding, hidden wealth, extrapolation or arbitrary equilibrium selection can yield uninformative sets.

\textbf{Causal interventions.} A potential outcome \(Y_i(p)\) is the outcome under a complete policy regime \(p\), including timing, financing, information and market spillovers. The target population and units are fixed before treatment. With interference, use the complete assignment vector or a justified exposure mapping; individual no-interference notation alone cannot identify an economy-wide intervention. A difference of mean potential outcomes does not require the cross-world joint distribution. Conditioning on post-treatment migration, survival, enrollment or employment changes the estimand and needs separate justification.

\textbf{Statistical statements.} An estimator, confidence set, forecast or test specifies a sampling law, model class and loss/coverage criterion. Uniform \((1-\alpha)\)-coverage means \(\inf_{m\in\mathcal M}\Pr_m\{\tau(m)\in C_n\}\geq1-\alpha\), or an explicitly asymptotic version. The whole parameter space trivially has coverage, so informativeness requires a stated width, risk or power objective. A forecasting improvement is relative to a named benchmark, information set and evaluation population.

\textbf{Equilibrium and optimization.} An equilibrium comprises feasible plans, optimality, beliefs where needed, budget constraints and all market/resource clearing. The primitives or a stated benchmark define each correspondence \(\mathcal E(p;m)\). Nonemptiness is assumed or proved, never inferred just from compact policy sets. A selected-equilibrium target uses a declared selection rule; a robust target quantifies over all permitted equilibria. Use suprema or infima unless attainment is established. An \(\varepsilon\)-optimal policy is feasible and within \(\varepsilon\) of the finite optimum value. Report an empty feasible set separately.

\textbf{Welfare and accounting.} Normative utility, interpersonal normalization, social weights, numeraire, discounting and the use of public receipts are primitives. Behavioral utility need not coincide with the welfare criterion. Household welfare includes ownership income and rents once; adding profits again to compensating variation generally double-counts them. A monetary equivalent is defined by an explicit transfer and price convention, with monotonicity and range conditions. Average effects, mechanism contributions, transfers and welfare losses are different targets. Infinite sums require convergence and dynamic feasible plans require terminal or no-Ponzi conditions.

\textbf{Source versions.} A working-paper date, revision date and journal publication year are distinguished. A DOI for a journal version is not automatically the DOI of an earlier manuscript. Locators refer to the stated version and printed pages where specified. A metadata-only check does not verify a claim about a full-text conclusion. Every entry's source subsection narrows the provenance claim accordingly.
% END ECONOMICS STATEMENT
\end{document}
